\subsection{What the global period estimate must control}
\label{shared-period-comparison-limits}

The period normalization distinguishes a large main contribution from
the positive trace error that enters the slope comparison. The following
results explain both the necessity of this distinction and the limits
of the pairwise estimates. All examples below concern the stated matrix
hypotheses; no realization by the theta construction is asserted.

Let $T_i:E_i\to H$ be finitely many complex linear maps, and put
\[
 R=\sum_i\dim E_i,\quad D=\dim H,\quad
 S(v_i)=\sum_iT_iv_i,\quad K=S^*S,\quad Q=SS^*,\quad M=\operatorname{tr}Q>0.
\]
Assume $D>0$. Write
\[
 \gamma=\frac1M\operatorname{tr}
   \left(\frac MR I_R-\bigoplus_iT_i^*T_i\right)_+,\qquad
 \beta=\frac1M\operatorname{tr}\left(Q-\frac MD I_D\right)_+.
\]
For the period maps the prescribed normalization is $T_i=c_i^{-1/2}P_i$,
with each $c_i>0$. When the diagonal identity contributions sum to $I_R$,
write $K=I_R+N$, retaining every remaining orbital contribution in $N$.

\begin{proposition}[The nonidentity sum is large]
\label{shared-period-redundancy}
If $r=\operatorname{rank}S$, then $N$ has $R-r$ eigenvalues equal to $-1$,
and
\[
 \operatorname{tr}N_-\ge R-r,\qquad
 \operatorname{tr}N_+\ge\max(0,M-r),\qquad
 \|N\|_{\rm op}\ge\max(0,M/D-1).
\]
Suppose, on these same full maps, that $0<\sigma\le1/20$ and
\[
 R\ge D_*/\sigma,\qquad D\le(1+\sigma/2)D_*,
 \qquad \sum_i\|T_i^*T_i-I\|_{\rm HS}^2\le e_0R,
 \qquad\sqrt{e_0}\le\sigma/2.
\]
Then
\[
 \frac{\operatorname{tr}N_+}{M}\ge1-2\sigma,
                 \qquad \|N\|_{\rm op}\ge\sigma^{-1}-2.
\]
\end{proposition}

\begin{proof}
The nonzero eigenvalues of $K$ are positive numbers $\lambda_j$ with
sum $M$. The eigenvalues of $N$ are $\lambda_j-1$ and $R-r$ copies
of $-1$. This proves the first assertions. Cauchy--Schwarz gives
$|M-R|\le\sqrt{R\sum_i\|T_i^*T_i-I\|_{\rm HS}^2}\le\sigma R/2$.
Hence $D/M\le\sigma(1+\sigma/2)/(1-\sigma/2)\le2\sigma$ and
$M/D-1\ge\sigma^{-1}-2$, proving the last assertions.
\end{proof}

The actual positive trace criterion already has the necessary shift:
\[
 \beta=\frac1M\operatorname{tr}
             \left(N-(M/D-1)I_R\right)_+.
\]
This follows from the equality of the positive spectra of $K$ and $Q$.
An absolute smallness estimate for $N$ would therefore impose an
inappropriate normalization. A possible reproducing identity must instead
concern $A=(D/M)K$: indeed
$\operatorname{tr}A=D$, $\operatorname{rank}A\le D$, and
$A^2=A$ is equivalent to $Q=(M/D)I_D$. Such an identity for the chosen
period family remains an additional global input.

Changes of comparison basis below mean new identifications with standard
Euclidean spaces, while the arithmetic metrics remain fixed separately.
Transporting the physical Gram matrices with the coordinates would leave
the spectral quantities unchanged and would give no optimization.

\begin{proposition}[Sharp trace loss in every comparison basis]
\label{shared-period-invariant-sharpness}
\label{shared-period-chirp}
\label{shared-period-simplex}
For each integer $m\ge1$, put $a=m^2+1$ and choose a prime $q\ge a^2$.
There are $q$ integral coordinate embeddings of common source rank
$r=2aq(q-1)$ into dimension $D=2q^2(q-1)$ such that
\[
 R=aD,\qquad T_i^*T_i=I_r,\qquad
 \|T_i^*T_j\|_{\rm HS}^2\le r^2/D\quad(i\ne j),
\]
but, over all independent invertible source bases and one common
invertible target basis,
\[
 \inf(\gamma+\beta)=\frac{m}{2a}
       =\frac12\sqrt{D/R}\sqrt{1-D/R}.
\]
The individual rank ratio $r/D=a/q$ can be arbitrarily small at fixed $m$.
\end{proposition}

\begin{proof}
Set $k_\pm=a\pm m$, and choose subsets $A_\pm\subset\mathbf F_q$ of
these cardinalities. Let $\mathcal G$ be the affine permutations
$z\mapsto uz+v$ of $\mathbf F_q$, and put $L=|\mathcal G|=q(q-1)$.
Index a target coordinate by $(g,\epsilon,z)\in
\mathcal G\times\{+,-\}\times\mathbf F_q$, and a source coordinate
by $(g,\epsilon,s)$ with $s\in A_\epsilon$. For $i\in\mathbf F_q$, define
\[
 T_i f_{g,\epsilon,s}=e_{g,\epsilon,s+g(i)}.
\]
These are isometric integral embeddings with $r=2aL$, $D=2qL$.
For $i\ne j$, the map $g\mapsto(g(i),g(j))$ runs once through the
ordered distinct pairs. A subset $A$ of size $k$ satisfies
$\sum_{b\ne c}|(A+b)\cap(A+c)|=qk(k-1)$, by counting at each
intersection coordinate. Therefore
\[
 \|T_i^*T_j\|_{\rm HS}^2
 =q\{k_+(k_+-1)+k_-(k_--1)\}=2q(a^2-1)
 \le2a^2(q-1)=r^2/D.
\]
Let $W$ be the plus-coordinate half of the target. Each plus coordinate
occurs in $k_+$ images and each minus coordinate in $k_-$ images, so
\[
 Q=k_+I_W\oplus k_-I_{W^\perp},\qquad M=R,\qquad
 \gamma=0,\qquad\beta=m/(2a).
\]
For the converse in every basis, put $V_i=T_i^{-1}W$.
Their dimensions are $Lk_+$, whereas $\dim W=D/2$.
For any family of maps and subspaces satisfying $T_iV_i\subset W$,
orthogonal projection onto these subspaces gives
\[
 \frac MR\sum_i\dim V_i-M\gamma
 \le\sum_i\|T_iP_{V_i}\|_{\rm HS}^2
 \le\operatorname{tr}(P_WQ)
 \le\frac MD\dim W+M\beta.
\]
Under $T_i\mapsto BT_iA_i$, use $A_i^{-1}V_i$ and $BW$; their
dimensions are unchanged. Hence $\gamma+\beta\ge
k_+/(2a)-1/2=m/(2a)$ in every comparison basis, attaining the claimed
infimum in the original coordinate bases.
\end{proof}

For any isometric family with the displayed pairwise bound,
\[
 \operatorname{tr}Q^2\le R+\frac{R^2-\sum_i r_i^2}{D},\qquad
 \beta\le\frac12\sqrt{\frac DR-\frac{\sum_i r_i^2}{R^2}}.
\]
Indeed, the centered target operator has trace zero, so its positive
trace is half its trace norm; apply Hilbert--Schmidt Cauchy--Schwarz.
The proposition asymptotically attains the constant $1/2$, even after
all comparison-basis optimization. In particular the pairwise contract
cannot imply a uniform bound $\inf(\gamma+\beta)=O(D/R)$.
Making the target sum equal to the identity by $T_i\mapsto Q^{-1/2}T_i$
merely transfers the defect: it gives $\beta'=0$ and
$\gamma'=m/(2a)$.

The example respects fixed arithmetic metrics as well as integrality.
Over $\mathbf Q$, give plus coordinate vectors norm $e^{-1}$ and minus
coordinate vectors norm $1$, with the induced source norms. The maps
have height zero, all source mean slopes are $p=k_+/(2a)$, and the
target mean slope is $1/2$. The source maximum and target minimum
slopes are $1$ and $0$. Thus their slope loss is exactly
$p-1/2=m/(2a)$. Auxiliary comparison-basis changes alter none of these
arithmetic quantities. A stronger bound for the actual period family
must therefore use additional information about those maps and the
subspaces they carry into a common target. No realization of these
examples by the theta construction is asserted.

